\documentclass[a4paper,11pt]{article}
\newcommand\subdir{/home/koen/LaTeX-setup/}
\input{\subdir/LaTeX-files/imports.tex}
\input{\subdir/LaTeX-files/master-internship.tex}
\usepackage{\subdir/LaTeX-files/aastex}
\usepackage{natbib}
\bibliographystyle{\subdir/LaTeX-files/astroadstitle.bst}

%Set the title, Author and date
\title{Notes Week 38}
\author{Koen Essers}
\tutorial{Master Internship}
\date{\today}

%Set the header style
\header{\tutorialstring}

%Begin the document
\begin{document}
\setuplayout
\section{Test}

My aim is to construct a way to predict the final period of a final ``simple'' binary calculation
based \emph{only} on the detailed single star model.

The easiest and most straightforward way to compute the final period is to use the simple binary
to trace the orbital evolution up to the point where the radius of the star is larger than the Roche lobe
radius.

\begin{figure}[ht]
	\centering
	\incplot{w38-direct}
\end{figure}

Now, let me see the final period for all of the models and compare them with detailed MESA models.

\begin{figure}[ht]
	\centering
	\incplot{w38-fracs-5}
\end{figure}

This is not great.
We see quite large differences probably largely due to the negligence of wind
during or after RLOF.

However, it seems quite difficult to, in advance, predict how much wind happens during and after RLOF.

\begin{figure}[ht]
	\marginfignote{0}{Here I show the fraction of the mass loss rate that is caused by wind versus the total
		mass loss rate.}
	\marginfignote{6}{The dots denote the rough start of RLOF.
		This is at the start of the most important
		``dip'' in this fraction.
	}
	\marginfignote{12}{Clearly, later models never really reach fractions close to 0, meaning winds will always
		remain important.}
	\centering
	\incplot{w38-fracs}
\end{figure}

\begin{figure}[ht]
	\marginfignote{0}{Here, I plot the change in mass due to winds and the total mass loss from \(q=0\) to some
		value \(q\).
		Here we can again clearly see that the fraction drops below unity once RLOF starts.
	}
	\marginfignote{9}{Now we again see that for the larger initial models, the fraction does not
		decrease to 0.}
	\centering
	\incplot{w38-fracs-cum}
\end{figure}

\begin{figure}[ht]
	\centering
	\incplot{w38-fracs-2}
\end{figure}

\bibliography{master.bib}
\end{document}
